\chapter{The SCIENCE of ALTERNATIVE DIMENSIONS}\label{ch:dimensions}
\markboth{The Science of Alternative Dimensions}{The Science of Alternative Dimensions}

\chapquote{``The greatest question of all is whether our experience on this planet is `it' or whether there is something else. Things in the supernatural realm give support, strangely perhaps, to the things we take on faith.''}{Art Bell}

\portrait[l]{0.30}{uf_kaku.jpg}{\textbf{\textcolor{ufogreenink}{Dr.~Michio Kaku}} (b.~1947), co-founder of string field theory and the man who taught more people what a dimension actually is than every physics department on Earth combined. Photograph: Cristiano Sant'Anna, CC0 / public domain.}{https://commons.wikimedia.org/wiki/File:Michio_Kaku-cropped.jpg}
This chapter belongs to one man, and unlike Chapter~\ref{ch:religion} the academy has no grounds to
object, because he is one of them. Michio Kaku holds the Henry Semat Chair in theoretical physics
at the City College of New York, he has held it for decades, and in 1974, with Keiji Kikkawa, he
wrote the first papers on \term{string field theory} --- the formalism that takes the strings of
string theory and puts them into a proper field-theoretic language, the way Maxwell's equations put
light into one. Every popular account of eleven dimensions you have ever heard, including the ones
misquoted at you by people claiming a craft came from one, traces back through his textbooks,
his lectures and his broadcasts. He is also the reason I am writing this book, so I am going to
spend a few pages on him and then use him as the instrument for the rest of the chapter.

He was born on 24 January 1947 in San Jose, California, and the first fact about him is not physics.
His parents were Japanese Americans who spent the war behind wire at the Tule Lake internment camp
in northern California --- which is where they met. His grandfather had come to America to help
clear the rubble of the 1906 San Francisco earthquake. So the man who ended up explaining the
structure of spacetime to the general public began as the child of two people the United States
had imprisoned for their ancestry, and he has never made a mystery of it or a grievance of it.

The second fact is the one that ought to be taught in every school. As a boy he read that Albert
Einstein had died with his great work unfinished, and he saw a photograph of the desk with the
incomplete unified field theory still on it. Most children who encounter that story feel a
pleasant sadness. Kaku decided he would finish it. Then --- and this is the part that makes
physicists laugh and parents go pale --- he decided that in order to work on the problem he needed
a particle accelerator, so in high school he built one. Not a model. A 2.3~MeV \term{betatron},
assembled in his parents' garage from about four hundred pounds of scrap transformer steel and
twenty-two miles of copper wire, which he wound by hand over the Christmas holidays on his high
school's football field because it was the only flat space long enough to lay the coils out.
The finished machine produced a magnetic field roughly twenty thousand times the Earth's and
successfully accelerated electrons, which is to say it worked. It also drew so much current that
it blew every fuse in the house. \hi{The man who would spend his career explaining energy scales
to the public learned about them the honest way: by tripping the domestic mains with a homemade
atom smasher.}

He took it to the National Science Fair in Albuquerque, where it caught the attention of Edward
Teller --- father of the hydrogen bomb, and not a man who handed out compliments --- and Teller
arranged a Hertz Engineering Scholarship that sent him to Harvard. He graduated summa cum laude
in 1968, first in his physics class, was drafted into the army during Vietnam and served two years
as an infantry trainee without being deployed, then took his doctorate at Berkeley in 1972 under
Stanley Mandelstam, went to Princeton on a lectureship, and by 1974 was publishing the string
field theory papers. \hi{Everything the popular press calls him --- futurist, television host,
explainer --- came after four decades of doing the actual mathematics.}

What he then did with that credential is the part relevant to this book. In \emph{Physics of the
Impossible} (2008) he took the entire toy box of science fiction --- force fields, invisibility,
teleportation, telepathy, faster-than-light travel, time machines, parallel universes,
precognition --- and sorted it into three classes: things impossible today but consistent with
known physics, things at the very edge of what the laws might permit, and things that would
require the laws themselves to be wrong. That taxonomy is the single most useful tool anyone has
ever handed ufology, and Chapter~\ref{ch:ancient} used its logic before I told you where it came
from. He wrote \emph{Physics of the Future}, \emph{The Future of the Mind} and \emph{The God
Equation} (2021), took the Klopsteg Memorial Award in 2008 and the Sir Arthur Clarke Lifetime
Achievement Award in 2021, presented the BBC's \emph{Time} in 2006 and \emph{Sci Fi Science} after
it, campaigned publicly against nuclear weapons with Peace Action, hosted a science programme on
WBAI for years, and --- because a man is allowed a hobby --- learned to figure skate.

He was also, repeatedly, a guest of the man quoted at the head of this chapter. Kaku went on
\emph{Coast to Coast AM} with Art Bell and told a late-night audience of millions that
extraterrestrial life is, on the numbers, a certainty rather than a speculation --- while declining
to say that any of it has landed here. That is exactly the line this book is trying to walk, and
he walked it on the paranormal radio show, in front of the most credulous audience available,
without once flattering them. \hi{You can be generous about the question and merciless about the
evidence. He proved it live on air.}

So: seven sections on the physics ufology most likes to borrow, each one measured against what the
man who built the mathematics actually says about it. Where he is optimistic, I will say so. Where
he has been accused of being too optimistic for public consumption, I will say that too, because
he would.

% ---------------------------------------------------------------
\section{11-Dimensional Spacetime}\label{sec:11dimensionalspacetime}

``Another dimension'' is the most abused phrase in ufology, and the abuse comes from a genuine
ambiguity in the English word. In physics a dimension is a degree of freedom --- an independent
number you need to specify a position. Our spacetime has four: three of space, one of time. In
popular usage a dimension is a \emph{place} --- a parallel world you could travel to. \hi{These are
not the same concept, and almost every interdimensional UFO claim silently swaps the first for
the second.}

\figwrap[r]{0.38}{A two-dimensional section of a Calabi--Yau manifold, the geometry proposed for the compactified dimensions.}{https://commons.wikimedia.org/wiki/File:Calabi-Yau-alternate.png}{fig:calabiyau}
The eleven-dimensional figure people cite is real and comes from a specific place. String theory
replaces point particles with one-dimensional vibrating strings, and its mathematics only stays
consistent in a particular number of dimensions --- ten for superstring theories, twenty-six for
the original bosonic version. In 1995 Edward Witten argued that the five known superstring
theories were limits of a single underlying eleven-dimensional framework, which he called
M-theory. The extra seven dimensions are proposed to be compactified: curled up at a scale near
the Planck length, around $10^{-35}$ metres, in configurations called Calabi--Yau manifolds (see Figure~\ref{fig:calabiyau}).

The standard analogy is a garden hose seen from a distance. It looks like a line --- one
dimension --- but an ant on its surface has two directions available, one of which wraps around
and is invisible from far away. Compactified dimensions are supposed to be like that wrapped
direction: everywhere, but too small to notice.

The garden-hose picture is Kaku's stock in trade, and it is worth noting how carefully he states
it when he is being precise rather than televisual. His formulation is that hyperspace is a
\emph{mathematical} necessity --- the equations of the fundamental forces become simpler and unify
when written in higher dimensions, and gravity in particular falls out of the geometry rather than
being bolted on. That is a claim about elegance and consistency. It is not a claim that the extra
dimensions are a corridor. \hi{Kaku's higher dimensions explain why the forces look the way they
do; they are not a place where anything keeps its spacecraft.}

Here is what matters for our purposes, and I am going to be blunt about it. If the extra
dimensions are curled up at $10^{-35}$ metres, then nothing can travel through them, because
there is nowhere to go. A nucleus is roughly $10^{20}$ times larger than that scale. Saying a
craft came from the eleventh dimension is like saying a lorry arrived via the gap between two
atoms. The claim isn't bold; it's dimensionally illiterate.

And string theory itself has a problem that ufologists ought to find sobering rather than
encouraging. It has, after fifty years, produced no confirmed experimental prediction. The
number of possible compactifications --- the ``landscape'' --- is estimated at something like
$10^{500}$, which means the framework can accommodate almost any low-energy physics and
therefore forbids almost nothing. A theory that cannot be wrong is in exactly the position
Section~\ref{sec:allevidencebutnoproof} described. Some very serious physicists, Lee Smolin and Peter Woit among them, argue
it has stopped being physics. Others expect it to be vindicated. The dispute is live, public,
conducted in journals, and settled by nobody yet --- which is what an honest hard problem looks
like.

Kaku is on the optimistic side of that argument and has been for fifty years, which is his right
as one of the people who built the thing. He calls the framework ``the theory of everything'' in
book titles and expects the Large Hadron Collider's successors, or cosmological signatures in the
polarisation of the microwave background, to eventually produce the missing evidence. His critics
--- Woit and Smolin most publicly --- say that language oversells a mathematical programme that has
yet to make a falsifiable prediction, and some of his fellow physicists have said the same about
his broadcasting more generally. I think both things are true at once, and it costs nothing to
say so: the field is unfinished, and the man who tells the public it is beautiful is not lying to
them about the mathematics. \hi{What he never does --- not once, in any book or broadcast --- is
tell you that anybody has travelled through a compactified dimension.} That step is always taken
by someone quoting him.

\begin{funfact}
Large extra dimensions are actually testable, and were tested. If some extra dimension were
millimetre-scale rather than Planck-scale, gravity would deviate from the inverse-square law at
short range, and energy would leak into it in high-energy collisions --- producing missing
energy, or even microscopic black holes, at the Large Hadron Collider. Physicists looked
carefully. Torsion-balance experiments have confirmed inverse-square gravity down to tens of
micrometres, and the LHC found no missing-energy signature. That is how you do it: take the
exotic idea, derive a consequence, go and look. Nobody found anything, and the search itself was
the achievement.
\end{funfact}

% ---------------------------------------------------------------
\section{Many Worlds Hypothesis}\label{sec:manyworldshypothesis}

The Many Worlds Interpretation is genuine physics and it is genuinely strange, so it is worth
getting right rather than gesturing at.

Quantum mechanics describes a system by a wavefunction which evolves smoothly and
deterministically according to the Schr\"odinger equation, and which generally represents a
superposition of possible outcomes. Measurement, however, yields a single definite result, with
probabilities given by the Born rule. Reconciling those two facts is the measurement problem,
and it is unsolved --- not in the sense that nobody has an answer, but in the sense that several
answers exist and no experiment currently distinguishes them.

The Copenhagen interpretation says the wavefunction collapses on measurement. This works and is
what most physicists use in practice, but it never satisfactorily defines what counts as a
measurement.

Hugh Everett, in his 1957 Princeton thesis, proposed removing collapse entirely. If the
Schr\"odinger equation always applies, with no exceptions for observers, then measurement
entangles the observer with the system and the superposition persists --- the observer simply
becomes a superposition of observers, each perceiving one definite outcome and each unaware of
the others. Decoherence, worked out later by Zeh, Zurek and others, explains why the branches
stop interfering: interaction with the environment scrambles the relative phases irreversibly and
fast.

The result is a picture with no collapse, no randomness, and no special role for observers ---
bought at the price of an unimaginable number of coexisting branches. Everett's work was ignored
for a decade, he left academia for defence analysis, and he died at fifty-one without seeing it
taken seriously.

Now the part relevant to this chapter. The branches of the MWI are, by the theory's own
mechanics, \emph{not accessible}. Decoherence is what makes branching effective, and it is
effectively irreversible for any macroscopic system --- the phase information is dispersed into
astronomically many environmental degrees of freedom. \hi{There is no mechanism in the formalism for
travelling between branches, sending a signal between them, or receiving a visitor from one.}
Many Worlds is not a multiverse of destinations. It is a claim about why measurement looks
probabilistic.

So when an interdimensional hypothesis borrows MWI for credibility, it is borrowing the name of
a theory whose central technical feature --- decoherence --- specifically forbids what is being
proposed. That is worse than having no theory. It is having one that says no.

This is where Kaku is most often misquoted, and it is worth being exact. He does discuss parallel
universes at length, he does think the multiverse follows naturally from inflationary cosmology
and from the string landscape, and in \emph{Physics of the Impossible} he classifies contact
between universes as a Class~II impossibility --- not forbidden by known law, but requiring
energies at the Planck scale, meaning something on the order of a stellar or galactic power
budget concentrated in a region smaller than a proton. \hi{``Not forbidden'' is doing an enormous
amount of work in that sentence, and the people who quote it always drop the energy figure.}
A civilisation able to do it would be a Type~III civilisation on the Kardashev scale, harnessing
the output of an entire galaxy. Kaku's own remark on such beings is that they would regard us
roughly as we regard the ants on a pavement --- which, if you take it seriously, is an argument
\emph{against} the interdimensional visitation narrative rather than for it. Nobody detours across
the multiverse to mutilate a cow.

\subsection{The Man from Taured}

Which brings me to the story that gets told in every lecture hall where this subject is discussed,
usually by someone who heard it from someone. In July 1954, the account goes, a well-dressed
European man arrived at Tokyo's Haneda Airport and presented a passport issued by the nation of
Taured. He gave his profession as business, said he had flown this route several times, and
produced a driving licence, a chequebook and other documents from the same country. Taured does
not exist. Shown a map, he indicated the territory of Andorra and became agitated, insisting his
country had existed for a thousand years. He was detained in a hotel room on an upper floor with
guards outside. The next morning the room was empty --- and so, reportedly, was the property
room holding his documents.

\figwrap[l]{0.38}{Tokyo's Haneda international arrivals hall in the 1950s --- the setting of a story with no surviving paperwork.}{https://commons.wikimedia.org/wiki/File:Haneda_Airport_1955.jpg}{fig:haneda}
It is a magnificent story. Now let me do to it what I have been asking you to do all term.
There is no contemporaneous Japanese press report (see Figure~\ref{fig:haneda}), no immigration record, no named officer, and
no airport file. The trail runs back to a 1981 anthology of strange stories by Colin Wilson and
John Grant, which retold it without a primary source, and behind that to a French tale of a
mysterious traveller circulating in the 1950s. There are two versions of the year, two of the
airport, and two of the country's name. \hi{Every element that would make the account checkable
is precisely the element that is missing.} What we have is not a case. It is a modern legend with
an unusually good hook, and its persistence is a fact about narrative appeal rather than about
cosmology --- which is worth noticing, because if the branches of the wavefunction were sending
us visitors, the evidence would look like paperwork, not folklore.

% ---------------------------------------------------------------
\section{The Mandela Effect}\label{sec:mandelaeffect}

A large number of people share confident, detailed, identical false memories. That is the actual
phenomenon, it is real, it is well documented, and the explanation is more interesting than the
parallel-universe reading.

The name comes from Fiona Broome, who around 2009 found many people who remembered Nelson
Mandela dying in prison in the 1980s. He was released in 1990 and died in 2013. The standard
catalogue: the Berenstain Bears remembered as Berenstein; ``Luke, I am your father'' (the line
is ``No, I am your father''); ``Mirror, mirror on the wall'' (``Magic mirror on the wall'');
the Monopoly man's monocle, which he does not have; Pikachu's tail tip, which is not black;
Curious George's tail, which he does not have; ``Life \emph{is} like a box of chocolates'' (it is
``\emph{was}''); and the Fruit of the Loom cornucopia, which never existed.

\hi{The mechanisms are established experimental psychology, and there are four of them working
together.}

\textbf{Memory is reconstructive.} This is the finding of Section~\ref{sec:confirmationbias} restated: recall is not
retrieval of a stored record but reconstruction from fragments, and each reconstruction can be
re-encoded with alterations. Elizabeth Loftus's work demonstrates this repeatedly and
quantitatively.

\textbf{Schema-driven completion.} We fill gaps with what fits. ``-stein'' is a far commoner
English surname ending than ``-stain'', so the brain regularises it. A wealthy cartoon
capitalist with a top hat and cane should have a monocle --- and Mr.~Peanut, who does, was in the
same visual culture. A cornucopia belongs on a fruit logo.

\textbf{Source confusion.} People remember content and lose its origin. Parodies, imitations,
misquotations in later media, and packaging redesigns all leave traces that get attributed to the
original. ``Luke, I am your father'' has been quoted that way in advertising and comedy for
forty years; the misquotation is more available than the line.

\textbf{Social contagion.} Once articulated, a false memory spreads and gains confidence.
Asking ``do you remember the cornucopia?'' plants it, and the internet asks at enormous scale.
And this is why the effect \emph{feels} paranormal: the shared quality is the surprising part,
but shared cognitive architecture plus shared media environment predicts shared errors. We all
run the same hardware on the same inputs. Of course we make the same mistakes.

Note what has quietly happened here. The Mandela Effect is habitually explained as evidence that
we have slipped between the branches of Section~\ref{sec:manyworldshypothesis} --- and the physics
of those branches, as its own architects describe it, forbids the slip. Kaku's energy requirement
for crossing between universes is Planck-scale. The energy requirement for misremembering a bear's
name is a cup of coffee and a badly designed logo. \hi{When two explanations differ by five
hundred orders of magnitude in cost, you are not choosing between theories. You are choosing
whether to think.}

\begin{funfact}
The Mandela Effect is the single most useful thing in this book, and I mean that. It is a
large-scale, self-administered demonstration that vivid, detailed, confidently held memory can be
simply and verifiably wrong --- because unlike a UFO sighting, you can check. The logo is
available. The film exists. Go and look, watch your own certainty survive the disconfirmation for
a second or two, and then remember that feeling next time you read a witness statement in
Chapter~\ref{ch:time}. Including mine.
\end{funfact}

% ---------------------------------------------------------------
\section{Life After Death}\label{sec:afterlife}

Near-death experiences are reported with striking cross-cultural consistency: a sense of leaving
the body, movement through darkness or a tunnel, light, encounters with figures or deceased
relatives, life review, a boundary, and reluctance to return. Bruce Greyson developed a
standardised scale for them, Raymond Moody popularised the topic in 1975, and the phenomenology
is real, well-catalogued, and often profoundly transformative for the person.

The neuroscience accounts for a surprising amount of it, feature by feature.

\textbf{Tunnel vision} is what cerebral hypoxia does to the visual system: peripheral retinal
and cortical function degrades before central, producing a constricting field with a bright
centre. Fighter pilots in centrifuge training experience it routinely and call it greyout,
along with euphoria and dreamlets during G-induced loss of consciousness --- James Whinnery's
research on this is directly relevant and largely ignored in NDE literature.

\textbf{Out-of-body experience} can be induced by stimulating the temporoparietal junction,
which integrates vestibular, proprioceptive and visual information into the sense of embodied
self. Olaf Blanke's work does this deliberately in the lab. Disrupt the integration and the
felt location of the self detaches. This is also what happens in sleep paralysis --- see
Section~\ref{sec:evidenceabductioncases} --- and in ketamine anaesthesia, which produces dissociation and NDE-like features
reliably enough to be studied.

\textbf{Life review} plausibly reflects disinhibited temporal lobe and hippocampal activity, and
temporal lobe seizures produce vivid autobiographical flooding along with intense religious
feeling.

\textbf{The 2013 Borjigin study} at Michigan found a surge of coherent high-frequency gamma
activity in rats after cardiac arrest, exceeding waking levels --- and comparable surges have
since been recorded in dying human patients, including a 2022 case captured incidentally during
EEG monitoring. A dying brain is not a quiet brain. It is briefly, intensely active.

The honest limits: the timing of NDE content relative to measurable brain activity is hard to
establish, and the AWARE studies led by Sam Parnia --- which placed images in resuscitation rooms
visible only from the ceiling, precisely to test veridical perception --- produced no confirmed
hits in the visual test across many years and hundreds of cardiac arrests. That is a genuine
negative result from a genuine attempt, and it deserves respect from both sides.

Why is this in a UFO book? Because Section~\ref{sec:closeencounters}'s abduction narrative and the NDE share their
entire structure: paralysis, a floating sensation, a passage, beings of light, a message about
humanity's future, a life review, and a return with transformed values. Two literatures
describing what is very plausibly one set of neurological events, interpreted through two
different cultural vocabularies. Whichever way that resolves, it tells us something.

Kaku's treatment of the mind is the useful frame here, and it is deliberately unromantic. In
\emph{The Future of the Mind} he defines consciousness as a process --- a brain building models of
itself in space and time in order to simulate the future --- and he treats memory, identity and
the felt sense of self as things a physical machine does, testable by interfering with the machine.
On that account an NDE is not a report from elsewhere but a report from a system under extreme
stress, which is precisely what the hypoxia, temporoparietal and gamma-surge findings above
describe. He is not dismissive about what comes next --- he talks seriously about connectomes,
memory recording and uploading as Class~I or Class~II problems --- but the direction of travel is
engineering, not the afterlife. \hi{The physicist whose entire career is about higher dimensions
looks at the tunnel of light and reaches for the oxygen curve, not the eleventh dimension.} That
ought to tell you something about how rarely the extra dimensions are actually needed.

% ---------------------------------------------------------------
\section{Shadow Creatures}\label{sec:shadow}

Shadow people --- dark humanoid figures seen peripherally, often at night, often accompanied by
dread, sometimes wearing a hat --- are reported consistently enough to constitute a real
perceptual phenomenon. The hat is the detail that convinces people, because how would strangers
independently invent a hat?

Several mechanisms produce this, and they stack.

\textbf{Hypnagogic and hypnopompic hallucination}, at sleep onset and waking, produces vivid
figures with full conviction of reality. Combined with sleep paralysis --- REM atonia persisting
into wakefulness --- it produces the complete package: immobility, a felt presence, a dark
intruder, crushing chest pressure, and terror. This is the Old Hag, the \emph{mara}, the
\emph{kanashibari}, the \emph{jinn} sitting on the chest, the incubus. Every culture on Earth has
a name for it, which tells you it is in the brainstem rather than the bedroom.

\textbf{Peripheral vision is rod-dominated}: high sensitivity, poor acuity, essentially no
colour. A peripheral figure is therefore reported as dark, featureless and human-shaped by
default, because that is the only kind of thing low-resolution monochrome peripheral processing
can return. The hat is likely the visual system resolving an ambiguous silhouette using the most
available human template.

\textbf{Pareidolia and face detection} run continuously and with a deliberate bias towards false
positives, because the cost asymmetry favours it. Mistaking a shadow for a person is cheap;
mistaking a person for a shadow was historically fatal.

\textbf{Felt presence} can be induced directly by stimulating the temporoparietal junction, and
is reported in high-altitude climbing, sensory deprivation, extreme fatigue, bereavement, and
solo endurance situations. Shackleton's fourth man; the ``third man factor''.

\textbf{Environmental contributors} are real and measurable: infrasound below about 20 Hz has been
associated with unease and shivers, and Vic Tandy's classic 1998 investigation traced a
``haunted'' laboratory to a 19 Hz fan resonance --- and to the resulting eyeball resonance
producing a grey figure at the edge of vision. Carbon monoxide poisoning produces dread and
hallucination and can kill you, which is why every account of a haunted bedroom should be met
with the question of whether anybody checked the boiler. Mould exposure and certain medications
contribute.

\begin{funfact}
The Hat Man deserves one more note, because the internet did something to him that is directly
relevant to Chapter~\ref{ch:abduction}. Before roughly 2000, shadow-figure reports were varied. After the hat
became a widely shared motif online, the hat became far more prevalent in new reports. That is
cultural transmission shaping the content of a genuine perceptual experience, in real time, in
an archive you can search. It is the same process as Section~\ref{sec:closeencounters}'s grey alien, and here we got to
watch it happen.
\end{funfact}

Kaku's taxonomy handles this class cleanly, and it is the one place in this chapter where he is
harsher than I am. Telepathy and telekinesis he rates Class~I --- not impossible at all, and in
fact already partly done: EEG and fMRI decoding, brain--machine interfaces, prosthetics driven by
motor cortex. What he denies is that any of it happens without hardware. A dark figure in a hat
at the edge of your vision is not a low-grade psychic event awaiting a physics; it is your retina
and your temporoparietal junction doing exactly what a lump of tissue does at three in the
morning. \hi{The honest physicist's position on shadow people is not that the paranormal is
impossible. It is that the machinery is already accounted for.}

% ---------------------------------------------------------------
\section{The Philadelphia Experiment}\label{sec:philadelphiaexperiment}

Before the case itself, a word about where we now are, because this section is the head of the
longest single thread in the book and it helps to know that in advance.
Section~\ref{sec:shadow} dealt with a figure manufactured inside one nervous system, from
ordinary perceptual machinery, with no author and no intent. What follows is the same
manufacture running at a larger scale --- on a document trail rather than a retina, with authors
and, eventually, with a production budget. The thread runs from here to Montauk
(Section~\ref{sec:montauk}), then through the Earthfiles cases of Chapter~\ref{ch:earthfiles} ---
mutilated cattle, hybrid animals, patterns in wheat, sounds from the sky, lights over a valley,
nine dead hikers, and animals that will not produce a body --- and it terminates on a
500-acre property in Utah (Section~\ref{sec:skinwalker}), which is where every element on that
list turns up at once. \hi{The thread is not a conspiracy joining the cases together. It is a
\emph{method} joining them, and by the last section you should be able to run it yourself
without my help.}

The story: in October 1943, at the Philadelphia Naval Shipyard, the destroyer escort USS Eldridge
was rendered invisible and teleported to Norfolk, Virginia and back, as part of a project to
defeat magnetic detection. Crew members were fused into the deck, driven mad, or vanished. The
technology derived from Einstein's unified field work and involved Nikola Tesla. It became a
1984 film, and it is one of the founding legends of the field.

The provenance is fully traceable, and this is one of the rare cases where we can watch a legend
assemble from parts.

\figwrap[r]{0.42}{USS \emph{Eldridge} (DE-173). Her deck log places her nowhere near the claimed events.}{https://commons.wikimedia.org/wiki/File:USS_Eldridge.jpg}{fig:eldridge}
It begins with Carl M.~Allen, writing as Carlos Miguel Allende, who in 1955--56 sent a series of
letters to the Office of Naval Research and to the astronomer Morris K.~Jessup, who had published
a book on UFOs. The letters were handwritten, in multiple ink colours, with erratic
capitalisation, and claimed Allende had witnessed the event from the deck of the SS Andrew Furuseth (see Figure~\ref{fig:eldridge}). He cited a Philadelphia newspaper story he could not name and never produced. He later
admitted to investigators that parts were invented. The annotated copy of Jessup's book that ONR
reproduced --- the ``Varo edition'' --- became scripture; ONR's own position, stated repeatedly
since, is that it never conducted any invisibility research.

The documentary record contradicts the story directly. The Eldridge's deck log survives. She was
commissioned in August 1943 and in October 1943 was conducting shakedown and convoy escort work
in the Atlantic; her movements are accounted for and she was not in Philadelphia on the claimed
date. Surviving crew members held reunions and denied it flatly, with several publicly annoyed
about it.

What the Navy \emph{was} doing at Philadelphia is genuinely interesting and is almost certainly
the seed. \textbf{Degaussing} --- wrapping a hull in energised cables to neutralise its magnetic
signature and defeat German magnetic mines and torpedoes --- was real, widespread, urgent, and
classified. It was routinely described in the yard as making a ship ``invisible'' to mines. That
is the whole trick: a real classified programme, a real piece of jargon, and a decade of
retelling.

The physics closes the case. Optical invisibility of a macroscopic steel object is not a matter
of insufficient power; metamaterial cloaking works only over narrow bandwidths at scales
appropriate to the wavelength, and a 100-metre destroyer across the visible spectrum is not a
harder version of the same problem. Teleporting 1,\!240 tonnes of steel and crew 200 miles has no
physical mechanism whatsoever, and Einstein never completed a unified field theory --- he worked
on it fruitlessly for thirty years, which is a matter of public record and a rather poignant one.

That last point is where this chapter's guide has personal standing, and I want to put it plainly.
The Philadelphia legend rests on the claim that Einstein's unified field theory was secretly
completed and secretly applied in 1943. A ten-year-old boy in California saw a photograph of the
desk where that work was left unfinished and spent the next sixty years trying to finish it ---
first with twenty-two miles of copper wire on a football field, then with string field theory, and
most recently in a book called \emph{The God Equation} whose entire premise is that the job is
\emph{still not done}. \hi{The unified field theory the Navy is supposed to have weaponised in
wartime Philadelphia is the same theory Michio Kaku is still working on, in public, in 2026.}
If it had been finished in a shipyard eighty years ago, the man who devoted his life to finishing
it would be the first to know, and he would have said so on the radio.

On invisibility itself Kaku is, characteristically, more generous than the legend deserves and
still fatal to it. He rates invisibility a Class~I impossibility --- achievable, in principle,
within known physics --- and points to metamaterials with negative refractive index as the
mechanism. But every working cloak so far operates over a narrow band of wavelengths at scales
matched to those wavelengths, which is the opposite of a 100-metre steel destroyer across the whole
visible spectrum. Teleporting that destroyer he places in the class where the physics does not
bend: quantum teleportation transfers states, not tonnage, and the state of $10^{30}$ atoms is not
a larger version of the state of one.

\begin{funfact}
Morris Jessup's involvement ended badly. He was a serious man --- he had done real astronomical
work on double stars --- and he was drawn into the Allende material, became increasingly isolated
and distressed, and died by suicide in 1959. That death was immediately absorbed into the legend
as evidence of silencing. It is the ugliest habit this field has: a man's collapse repurposed as
corroboration. Whatever else you take from this book, don't do that.
\end{funfact}

\section{The Montauk Project}\label{sec:montauk}

The Philadelphia Experiment has a sequel, and it is instructive because we can date its
invention almost to the month.

The story: after 1943, the Navy's invisibility work did not stop --- it went underground, quite
literally, beneath Montauk Air Force Station at the eastern tip of Long Island. There, from the
late 1960s to 1983, a project called Phoenix ran experiments in electromagnetic mind control,
weather modification, materialisation of objects out of thought, and finally time travel, using
a device called the Montauk Chair. On 12 August 1983 the Montauk apparatus locked onto the USS
Eldridge across a forty-year gap, closing a loop between the two experiments. The project
ended when a participant deliberately released a monster from his own subconscious, which
rampaged through the base and forced its abandonment.

I want to be fair about the appeal here. As mythology this is superb --- it has a haunted
military base, a chair that reads minds, a time loop, and a creature made of thought. As
history it has nothing at all, and the useful exercise is separating the real base from the
story bolted onto it.

\figwrap[r]{0.42}{The AN/FPS-35 radar at Camp Hero. The 126-foot reflector still stands; the
building beneath it is a decommissioned Cold War radar station, not a laboratory.}{https://commons.wikimedia.org/wiki/File:Montauk_Air_Force_Station_radar_tower.jpg}{fig:montaukradar}
Camp Hero is real, and genuinely strange to look at. Fortified in 1942 as Fort Hero, it mounted
16-inch coastal guns against a feared German naval attack, and its buildings were deliberately
disguised as a sleepy New England fishing village --- pitched roofs, dormer windows, a mock
church --- so that a submarine periscope would see nothing worth shelling. In December 1960 the
site brought online an AN/FPS-35 long-range search radar (see Figure~\ref{fig:montaukradar}): a
126-foot reflector weighing some forty tons, with a range beyond 200 miles, part of the
continental air-defence net watching for Soviet bombers. It interfered badly with local
television and radio reception, was taken off the air in 1961 and returned in 1962. The station
closed on 31 January 1981, and the radar was left in place --- not out of secrecy, but because
local fishermen had come to rely on the tower as a navigational landmark. The land became Camp
Hero State Park, opened to the public on 18 September 2002, with sections restricted for
unexploded ordnance from the gunnery era.

So the ingredients were sitting there in plain sight: a decaying secret-looking base, a giant
rotating antenna that demonstrably interfered with the electronics in people's homes, sealed
bunkers, and a fence.

\hi{The legend's provenance is not obscure, contested, or classified --- it was written by two
identifiable men and published by their own imprint.} Preston Nichols, an electronics engineer
born in 1946, began telling the story in the early 1980s, claiming he had recovered his own
suppressed memories of working at the base. With Peter Moon he published \emph{The Montauk
Project: Experiments in Time} in 1992, followed by a long series of sequels of escalating
strangeness. Al Bielek attached himself to the account, claiming under regression that he was
really ``Edward Cameron'', a crewman of the Eldridge who had jumped forward from 1943 --- a
claim undermined by his birth certificate, which has him born in 1927, and by the complete
absence of any Cameron in the Eldridge's records. Duncan Cameron supplied the psychic-chair
material. The dates were chosen for numerological reasons: 12 August 1983 was said to sit at
the peak of a twenty-year biorhythm cycle that resonated with August 1943.

Jacques Vallée, who is not a debunker by temperament and has spent his career arguing that the
phenomenon deserves serious study, examined the Montauk and Philadelphia material and published
his verdict in the \emph{Journal of Scientific Exploration} in 1994 under the title ``Anatomy of
a Hoax''. His conclusion was that the entire structure rests on recovered-memory testimony from
a handful of interlinked people, with no document, no payroll record, and no independent
witness. The authors themselves described their content as ``soft facts'' and wrote that the
reader could take it ``as science fiction or non-fiction''; libraries catalogue the series as
speculative fiction. That is not a hostile characterisation --- it is the authors' own framing.

The recovered-memory point is the whole case, and I will not repeat the argument here beyond
noting that everything in Section~\ref{sec:hypnosis} applies with full force: hypnotic
regression does not retrieve buried experience, it generates confident narrative, and a story
whose only evidence is regressed testimony has no evidentiary base at all. Montauk is
Section~\ref{sec:philadelphiaexperiment} run through that process a second time, forty years
later, with the added detail that this time we can name the authors and read their copyright
page.

What I want you to carry out of this section is not the verdict but the recipe, because you are
about to meet it repeatedly. A decommissioned installation and a real classified technology
supply the setting. Regressed testimony from a small interlinked group supplies the events.
A numerological coincidence supplies the link to an earlier legend, which lends borrowed
age and depth. Published sequels supply the corroboration, each citing the last. And a
television production eventually supplies the audience, at which point the story stops needing
evidence altogether because it has something better --- a schedule.

Hold that last item in mind. Chapter~\ref{ch:earthfiles} takes the same recipe out of the bunker
and into open country, where the ingredients are a dead animal, a flattened field, a noise, a
dead climbing party, and a missing ape; and Section~\ref{sec:skinwalker}, where this thread ends,
is the case where a property with a genuinely explicable geology acquired all of them at once
and then, precisely as Montauk did, acquired a camera crew.

\begin{funfact}
The Duffer brothers' original working title for \emph{Stranger Things} was ``Montauk'', and the
first season was written around a Long Island setting with a secret government lab, a missing
boy and a creature from another dimension before the production moved to Georgia and renamed
the town Hawkins. The legend also produced a 2015 documentary, \emph{Montauk Chronicles}. If
the Montauk Project has a lasting legacy, it is as set design --- which, given that its authors
told us to read it as fiction, is arguably exactly what they asked for.
\end{funfact}

And since Montauk's central prop is a time machine, it is worth closing the chapter with the
quotation at the head of Chapter~\ref{ch:time}, which belongs to the man whose portrait opens this
one. Kaku's line is that the burden of proof has shifted --- that physicists once had to show time
travel was possible and now have to show that some law forbids it. \hi{Read that carefully,
because it is an invitation to do mathematics, not an announcement that somebody already has.}
Wormholes, closed timelike curves and negative energy densities are live topics in general
relativity, argued in journals, and the boundary conditions are brutal: exotic matter nobody has
isolated, energies nobody can reach, and Hawking's chronology-protection conjecture waiting at the
end of it. That is the difference between the two halves of this chapter. One half is a hard
problem being worked on by people who publish their failures. The other is a chair in a bunker on
Long Island, remembered under hypnosis, catalogued by librarians as fiction.

\begin{funfact}[Fun Fact: Twenty-Two Miles of Copper Wire]
Go back to the football field for a moment, because the numbers deserve to be felt rather than
read. Twenty-two miles of copper wire is roughly the distance from Manhattan to Yonkers and back,
and a teenager wound it by hand, in coils, over a school holiday, on the fifty-yard line ---
because his parents' garage was not long enough to lay it out straight. Four hundred pounds of
scrap transformer steel came from a junkyard. The finished betatron reached 2.3~MeV and produced a
magnetic field about twenty thousand times stronger than the Earth's, and when he switched it on it
took out every fuse in the house. His mother's recorded reaction was to wonder why her son could
not find a nice job and marry a nice girl instead of building machines in the garage. He tells the
story himself in the introduction to \emph{Physics of the Impossible} --- it is his own account, not
a second-hand legend, which is exactly the standard this book keeps asking of everybody else. He
went to the Albuquerque science fair with it, Edward Teller noticed, a scholarship to Harvard
followed, and
fifty years later he was explaining eleven-dimensional spacetime to several million people at a
time. \hi{Every single item in that sequence is cheaper than a Planck-scale portal and none of it
required another dimension.} That is what actually doing the physics looks like: junkyard steel,
blown fuses, an angry electricity meter, and a lifetime of work on somebody else's unfinished
equation.
\end{funfact}

\begin{srcnotes}
\item Kaku, M.\ \emph{Hyperspace: A Scientific Odyssey Through Parallel Universes, Time Warps and
the Tenth Dimension} (Oxford University Press, 1994) --- the garden-hose analogy, compactification,
and higher-dimensional unification as a mathematical rather than geographical claim.
\item Kaku, M.\ \emph{Physics of the Impossible} (Doubleday, 2008) --- the Class~I/II/III taxonomy
used throughout this chapter, including invisibility and metamaterials, telepathy and telekinesis,
teleportation, parallel universes and time travel. This is also Kaku's own first-hand source for
the garage betatron: the Einstein desk photograph, the 2.3~MeV machine, the twenty-two miles of
hand-wound copper wire on the school football field and the blown fuses are all recounted by him in
the book's introduction.
\item Kaku, M.\ \emph{The Future of the Mind} (Doubleday, 2014) and \emph{The God Equation}
(Doubleday, 2021) --- consciousness as a physical modelling process; the unified field theory
explicitly described as still unfinished.
\item Kikkawa, K.\ \& Kaku, M.\ (1974). Field theory of relativistic strings.
\emph{Physical Review~D} --- the founding papers of string field theory.
\item Biographical details on Kaku --- birth 24 January 1947 in San Jose; his parents' internment
at Tule Lake and meeting there; his grandfather's arrival for the 1906 San Francisco earthquake
clean-up; the Albuquerque National Science Fair and Edward Teller's Hertz
scholarship; Harvard summa cum laude 1968; army service 1968--70; Berkeley PhD 1972 under Stanley
Mandelstam; the Henry Semat Chair at CCNY/CUNY; the Klopsteg Memorial Award 2008 and Sir Arthur
Clarke Lifetime Achievement Award 2021 --- from his own published accounts, his CCNY faculty
record, and his standard biographical entries.
\item Kaku's \emph{Coast to Coast AM} appearances with Art Bell, including the 30~November 2007
broadcast in which he described extraterrestrial life as a statistical certainty while declining
to endorse any landing claim.
\item Witten, E.\ (1995). String theory dynamics in various dimensions. \emph{Nuclear Physics~B}
--- the eleven-dimensional M-theory conjecture.
\item Woit, P.\ \emph{Not Even Wrong} (Jonathan Cape, 2006); Smolin, L.\ \emph{The Trouble with
Physics} (Houghton Mifflin, 2006) --- the falsifiability case against the string programme cited
above.
\item Everett, H.\ (1957). ``Relative state'' formulation of quantum mechanics.
\emph{Reviews of Modern Physics}~29, 454; with Zurek, W.\ H.\ (2003), Decoherence, einselection,
and the quantum origins of the classical, \emph{Reviews of Modern Physics}~75, 715.
\item Vall\'ee, J.\ (1994). Anatomy of a hoax: the Philadelphia Experiment fifty years later.
\emph{Journal of Scientific Exploration}~8(1); Nichols, P.\ \& Moon, P.\ \emph{The Montauk
Project: Experiments in Time} (Sky Books, 1992).
\item USS \emph{Eldridge} (DE-173) deck logs and Office of Naval Research statements on the
Philadelphia Experiment, US National Archives and ONR public information releases.
\item Borjigin, J.\ et al.\ (2013). Surge of neurophysiological coherence and connectivity in the
dying brain. \emph{PNAS}~110(35); Parnia, S.\ et al.\ (2014). AWARE---AWAreness during
REsuscitation. \emph{Resuscitation}~85(12); Blanke, O.\ et al.\ on induced out-of-body experience;
Whinnery, J.\ E.\ on G-induced loss of consciousness and dreamlets.
\item Tandy, V.\ \& Lawrence, T.\ (1998). The ghost in the machine. \emph{Journal of the Society
for Psychical Research}~62 --- the 19~Hz infrasound case.
\end{srcnotes}
